\documentclass[11pt]{article}
\usepackage[margin=1in]{geometry}
\usepackage{amsmath,amssymb,amsthm,booktabs}
\usepackage[hidelinks]{hyperref}
\setlength{\emergencystretch}{2em}
\newtheorem{theorem}{Theorem}
\newtheorem{lemma}[theorem]{Lemma}
\newtheorem{corollary}[theorem]{Corollary}
\newcommand{\mr}{\operatorname{mr}}
\newcommand{\rank}{\operatorname{rank}}
\title{A universal reflection obstruction for sixteen knots}
\author{Internal research working proof}
\date{17 November 2026}
\begin{document}
\maketitle

The sixteen alternating knots listed in Theorems~\ref{thm:alternating8}
and~\ref{thm:a6eight} have meridional rank and bridge number four by
retained finite-action certificates. Nevertheless, no Coxeter
reflection quotient, finite or infinite, can certify that sharp lower
bound. A short calculation with independent reflection directions
proves the obstruction over every field of characteristic different
from two. The form may be indefinite or singular. The argument uses
no Coxeter-group classification and no determinant localization.
The theorem is a written proof with exact internal algebra checks;
formal verification, external review and worldwide priority are not
claimed.

\section{The reflection statement and its exact diagram}

Let $k$ be a field of characteristic different from two, and let $B$
be a symmetric bilinear form on a $k$-vector space. For $B(u,u)=2$ put
\[
                         R_u(x)=x-B(u,x)u.
\]
This reflection is an involution preserving $B$. No nondegeneracy
assumption is made. Its displacement image is the line $ku$, since
$(R_u-I)u=-2u$. Consequently
\begin{equation}\label{eq:sign}
                   R_u=R_v\quad\Longleftrightarrow\quad v=\pm u
                   \qquad(B(u,u)=B(v,v)=2).
\end{equation}
Indeed equality forces their displacement lines to agree, and the
norm condition makes the proportionality scalar have square one.

Use the following retained planar diagram for $K=13a_{373}$:
\[
\begin{split}
 &(3,1,4,26),\ (1,7,2,6),\ (7,3,8,2),\ (13,4,14,5),\\
 &(5,14,6,15),\ (21,9,22,8),\ (9,17,10,16),\ (17,11,18,10),\\
 &(11,25,12,24),\ (19,12,20,13),\ (15,22,16,23),\\
 &(23,18,24,19),\ (25,21,26,20).
\end{split}
\]
An arc is a component of edge labels joined through the overpasses
(slots one and three, with slots numbered zero through three).
Sort these components lexicographically by their increasing edge lists
and number the arcs from zero. The four seed arcs are $1,2,6,8$.
A triple $(o,a,b)$ below means that the over-arc $o$ conjugates the
reflection on under-arc $a$ to the reflection on under-arc $b$.
Successive propagation uses
\begin{equation}\label{eq:trace}
\begin{gathered}
 (2,6,7),\ (7,2,3),\ (3,1,0),\ (1,3,4),\ (8,4,5),\\
 (5,8,9),\ (6,9,10),\ (10,0,12),\ (4,10,11).
\end{gathered}
\end{equation}
The four remaining crossing triples, up to swapping the two under-arcs,
are
\begin{equation}\label{eq:left}
                    (0,1,2),\quad(12,5,6),\quad
                    (11,7,8),\quad(9,11,12).
\end{equation}
The retained named-census transfer and planar check bind this literal
diagram to $13a_{373}$. The present checker reconstructs its arc
components and verifies every listed propagation and remaining crossing.
It does not infer knot identity from the name attached to a diagram.

\begin{theorem}\label{thm:reflection}
In every homomorphism from $\pi_1(S^3\setminus13a_{373})$ whose
meridians are nonisotropic orthogonal reflections preserving one
symmetric bilinear form over $k$, the normal directions of the four
specified seed meridians span a space of dimension at most three.
This holds for every field $k$ of characteristic different from two,
including singular forms and ambient spaces of any dimension.
\end{theorem}

Conjugacy of knot meridians allows their normal vectors to be chosen
with one common nonzero norm. Rescaling $B$ makes that norm two.
Thus the normalization used in the theorem requires no extraction of
square roots in $k$.

\section{The first loop reduces the Gram matrix}

Suppose, for a contradiction, that the four seed normal vectors are
independent, and use them as the coordinate basis
\[
                  e_0=u_1,\quad e_1=u_2,\quad
                  e_2=u_6,\quad e_3=u_8.
\]
Their span is invariant under all four reflections. Restricting $B$
to this span gives an arbitrary symmetric matrix with diagonal two;
write $a_{ij}=B(e_i,e_j)$ for its six off-diagonal entries.
Along \eqref{eq:trace}, choose each new normal vector to be exactly
$u_b=R_{u_o}(u_a)$. This is permitted by \eqref{eq:sign} and preserves
norm two. At the unused crossings the two normal vectors may differ
by either sign; no sign choices are discarded.

Put $t=a_{12}$. The first three propagation steps give
\[
 \begin{split}
 u_7&=e_2-te_1,\\
 u_3&=(1-t^2)e_1+te_2,\\
 u_0&=e_0-su_3,\qquad
             s=B(u_3,e_0)=(1-t^2)a_{01}+ta_{02}.
 \end{split}
\]
The first remaining crossing says $R_{u_0}(e_0)=\varepsilon e_1$
for $\varepsilon\in\{1,-1\}$. Set $H=B(u_0,e_0)$.
Comparison of coefficients in the independent vectors $e_0,e_1,e_2$
gives
\[
                  H=1,\qquad s(1-t^2)=\varepsilon,\qquad st=0.
\]
The second equation proves $s\ne0$ before any cancellation is used.
It follows that $t=0$ and $a_{01}=s=\varepsilon$.
Replace the initial choice of root $e_1$ by $\varepsilon e_1$.
The reflection on that seed is unchanged, all norms remain two,
and the recomputed trace normalizes
\[
              a_{01}=1,\qquad a_{12}=0,\qquad
              u_0=e_0-e_1,\qquad u_4=e_1-e_0.
\]
In particular, this reduction never divides by a Gram determinant or
assumes that $B$ has trivial radical.

\section{Two further crossings are inconsistent}

Put
\[
              U=e_0-e_1,\qquad V=e_2,\qquad W=e_3.
\]
These vectors are independent, because the original four seed vectors
are a basis. They all have norm two. Write
\[
 x=B(U,V),\qquad y=B(U,W),\qquad z=B(V,W),\qquad h=1-y^2.
\]
The remaining trace directions depend only on these three vectors.
Define
\[
                 v=xy+zh,\qquad w=3y-y^3-xv.
\]
Direct reflection gives
\begin{align*}
 u_5&=-U+yW,\\
 u_9&=yU+hW,\\
 u_{10}&=yU-vV+hW,\\
 u_{11}&=(y-w)U-vV+hW,\\
 u_{12}&=(1-wy)U+wvV-whW.
\end{align*}
Here $w=B(U,u_{10})$. These formulas use only symmetry and the
diagonal norms; no value or determinant of the three-dimensional
Gram matrix has been specified.

The unused crossing $(11,7,8)$ says
$R_{u_{11}}(V)=\beta W$ for $\beta\in\{1,-1\}$.
Putting $p=B(u_{11},V)$ and comparing the three independent
coefficients yields
\[
             p(y-w)=0,\qquad 1+pv=0,\qquad -ph=\beta.
\]
The middle equation proves $p\ne0$. Hence
\[
                              w=y,\qquad v=\beta h.
\]
The latter follows by comparing $pv=-1$ and $ph=-\beta$;
only the already proved nonzero scalar $p$ is cancelled.
Thus
\[
                         u_{12}=hU+\beta yhV-yhW.
\]
The unused crossing $(12,5,6)$ says
$R_{u_{12}}(u_5)=\alpha V$, where $\alpha\in\{1,-1\}$.
Set $q=B(u_{12},u_5)$. Its $U$-coefficient and $V$-coefficient give
\[
                         qh=-1,\qquad -q\beta yh=\alpha.
\]
Consequently $\beta y=\alpha$, so $y^2=1$ and $h=0$.
This contradicts $qh=-1$, proving Theorem~\ref{thm:reflection}.
The last remaining crossing $(9,11,12)$ is not needed.
No division by two, three or five occurs in this contradiction;
characteristic different from two is needed for the initial reflection
and sign normalization, rather than for a discarded denominator.

\section{Three sibling knots}

\begin{theorem}\label{thm:siblings}
The conclusion of Theorem~\ref{thm:reflection} also holds for
$13a_{375},13a_{391},13a_{392}$ and their original retained four-seed
systems. The field may be any field of characteristic different
from two, and the common form may be singular.
\end{theorem}
\begin{proof}
All diagrams and original seed systems are retained in the same input.
The local four-step trace
\[
                 (2,6,7),\ (7,2,3),\ (3,1,0),\ (1,3,4)
\]
and closing crossing $(0,1,2)$ occur in every case. To apply the
first-loop reduction with four independent normals, use the following
seed systems:
\[
\begin{array}{c|c|c}
 K&\text{original seeds}&\text{new seeds}\\\hline
13a_{375}&(1,2,6,9)&(1,2,6,9)\\
13a_{391}&(1,2,5,9)&(1,2,6,9)\\
13a_{392}&(1,2,5,11)&(1,2,6,12).
\end{array}
\]
These changes preserve independence, without a Gram determinant
hypothesis. For $13a_{391}$, if $P=u_5,Q=u_9$ and
$z_0=B(P,Q)$, the new pair is $(P-z_0Q,Q)$, of determinant one.
For $13a_{392}$, put $P=u_5,Q=u_{11}$ and $z_0=B(P,Q)$.
Then
\[
 u_{12}=Q-z_0P,\qquad
 u_6=(1-z_0^2)P+z_0Q.
\]
The transformation from $(P,Q)$ to $(u_6,u_{12})$ has determinant
$(1-z_0^2)+z_0^2=1$. All roots still have norm two.
Thus, under the supposition of independent original seed directions,
the first-loop proof again gives independent vectors
$U=u_0=-u_4$, $V=u_6=u_7$, and $W$, where $W=u_9$ in the first
two cases and $W=u_{12}$ in the third.

The additional five trace steps are, respectively,
\[
\begin{array}{c|l}
13a_{375}&(6,9,10),(10,7,8),(4,10,11),(11,4,5),(8,11,12)\\
13a_{391}&(9,6,5),(6,9,10),(4,10,11),(5,11,12),(11,7,8)\\
13a_{392}&(12,6,5),(5,12,11),(4,11,10),(6,10,9),(10,7,8).
\end{array}
\]
Every step is a literal crossing or its reversed under-arc direction.
In the remainder write $x=B(U,V)$, $y=B(U,W)$ and $z=B(V,W)$.

For $13a_{375}$, put $T=R_VW=W-zV$,
$A=R_UT=T-tU$, and $t=B(U,T)=y-zx$.
Then
\[
 u_8=R_TV,\quad u_{11}=A,\quad
 u_5=-R_AU=(t^2-1)U-tT,\quad u_{12}=R_{u_8}A.
\]
The vector $u_8$ has no $U$-coefficient, so $u_{12}$ has
$U$-coefficient $-t$. The unused crossing $(9,12,0)$ says
$R_Wu_{12}=\gamma U$, so
$u_{12}=\gamma R_WU=\gamma(U-yW)$ and $t=-\gamma$.
Hence $t^2=1$ and $u_5=-tT$. The unused crossing $(5,8,9)$
then gives
\[
                \alpha W=R_{u_5}u_8=R_TR_TV=V,
\]
contrary to independence of $V,W$.

For $13a_{391}$ use the same $T,A,t$ and put $D=R_WV=V-zW$.
The trace gives
\[
 u_5=D,\qquad u_{11}=A,\qquad
 u_{12}=R_DA=A-rD,\quad r=B(D,A),\qquad u_8=R_AV.
\]
The unused crossing $(10,12,0)$ says
$u_{12}=\gamma R_TU=\gamma(U-tT)$.
Comparison of $U$-coefficients gives $t=-\gamma$.
The two $V$-coefficients are then $-z-r$ and $-z$, so $r=0$.
Thus $u_{12}=A$. The unused crossing $(12,8,9)$ yields
\[
                      \beta W=R_{u_{12}}u_8=R_AR_AV=V,
\]
again impossible.

For $13a_{392}$ put $h=1-z^2$ and
$t=B(U,u_{11})$. The trace yields
\[
\begin{split}
 u_5&=V-zW,\qquad u_{11}=zV+hW,\\
 u_{10}&=-tU+zV+hW,\\
 u_9&=-tU+(z-k)V+hW,\qquad k=B(V,u_{10}),\\
 u_8&=ktU+(1-kz)V-khW.
\end{split}
\]
The unused crossing $(9,12,0)$ says $R_{u_9}W=\gamma U$.
Writing $p=B(u_9,W)$ and comparing coefficients gives
\[
                       pt=\gamma,\qquad p(z-k)=0,\qquad ph=1.
\]
Thus $p,t,h$ are nonzero and $k=z$. The unused crossing
$(11,8,9)$ says $R_{u_{11}}u_8=\beta u_9$.
Since $u_{11}$ has no $U$-coefficient, its $U$-coefficients give
$kt=-\beta t$. Hence $z=k=-\beta$, so $h=0$, contradicting $ph=1$.
All cancellations used a scalar already proved nonzero; no division
by a Gram determinant or characteristic-dependent constant occurs.
This proves the theorem.
\end{proof}

\section{The other four alternating seeds}

For each of the four cases already proved, identify arcs $0,4$ as one
reflection direction and arcs $6,7$ as another. Remove the five local
crossings involving arcs $1,2,3$. There are eight crossing relations
on the remaining directions labelled
\[
                           0,5,6,8,9,10,11,12.
\]
Here the identifications concern reflection lines; the chosen vectors
were $u_4=-u_0$ and $u_7=u_6$.

\begin{lemma}\label{lem:reduced}
In each of the four reduced eight-crossing systems just described,
the span of all its normal directions has dimension at most two.
\end{lemma}
\begin{proof}
The displayed traces for the four cases propagate every remaining
normal direction from $U=u_0,V=u_6$ and a third root:
$u_8,u_9,u_9,u_{12}$, respectively. Reflection of vectors in a linear
subspace stays in that subspace when its normal also belongs to it.
Thus these three normals span every reduced direction.
If they were independent, the three-coordinate contradictions in the
preceding two sections would apply verbatim. Those arguments used
only the eight reduced crossings and their norm-two reflection
formulas, rather than any unlisted entry of the original Gram matrix.
Consequently the three normals are dependent and their span has
dimension at most two.
\end{proof}

Put
\[
\mathcal K_8=\{13a_{373},13a_{375},13a_{391},13a_{392},
13a_{3264},13a_{3266},13a_{3275},13a_{3276}\}.
\]

\begin{theorem}\label{thm:alternating8}
For every knot in $\mathcal K_8$, the normal directions of the four
original retained seed meridians are dependent in any common-form
nonisotropic orthogonal reflection representation over a field of
characteristic different from two. The form may be singular.
\end{theorem}
\begin{proof}
The first four knots have already been proved. For the other four,
suppose that their original seed normals are independent. The
following changes give another normal basis:
\[
\begin{array}{c|c|c}
 K&\text{original seeds}&\text{changed seeds}\\\hline
13a_{3264}&(1,3,5,10)&(1,3,6,10)\\
13a_{3266}&(1,3,5,8)&(1,3,6,9)\\
13a_{3275}&(1,3,5,9)&(1,3,6,9)\\
13a_{3276}&(1,3,5,9)&(1,3,6,9).
\end{array}
\]
Here are their independence justifications before any Gram conditions
are imposed. Write $e_0=u_1,e_1=u_3$, and initially let $P,Q$ be
the last two original seed normals. For $13a_{3264}$ and
$13a_{3276}$, $u_0=R_{e_1}e_0$ lies in the first two-dimensional
plane, $u_{12}=R_Qu_0$ lies in $\langle e_0,e_1,Q\rangle$, and
$u_6=R_{u_{12}}P$ has coefficient one on $P$. Replacing $P$ by
$u_6$ therefore has determinant one.
For $13a_{3266}$ first set $u_9=R_PQ$, keeping $P$; this is a
unit-determinant change of the last pair. Now $u_{12}=R_{u_9}u_0$
and $u_6=R_{u_{12}}P$, so the same triangular argument applies.
For $13a_{3275}$ the change is simply $u_6=R_QP$, keeping $Q$.
Thus all changed seed lists are bases even if the Gram form is singular.

Set $V=u_6$ and let $W_0$ be the last root in the changed seed list.
The five literal crossings of the initial local region are
\[
              (3,1,0),\ (1,3,4),\ (0,3,2),\ (2,6,7),\ (7,1,2).
\]
Write $a=B(e_0,e_1)$. Choosing propagated normal vectors gives
\[
 u_0=e_0-ae_1,\quad u_4=e_1-ae_0,\quad
 u_2=ae_0+(1-a^2)e_1,\quad u_7=R_{u_2}V.
\]
The vector $u_7$ has coefficient one on $V$, since $u_2$ lies in
$\langle e_0,e_1\rangle$. The closing crossing says
$R_{u_7}e_0=\varepsilon u_2$. Its $V$-coefficient is
$-B(u_7,e_0)=0$, so $R_{u_7}e_0=e_0$ and
$u_2=\varepsilon e_0$. Coefficient comparison gives
$a=\varepsilon$ and $a^2=1$. Flip the sign of $e_1$ if necessary to
normalize $a=1$. This gives
\[
                  u_0=U=e_0-e_1,\quad u_4=-U,\quad u_2=e_0.
\]
Since $u_7=R_{e_0}V$ and $B(u_7,e_0)=0$, preservation of $B$ gives
$B(V,e_0)=0$ and $u_7=V$. The three vectors $U,V,W_0$ are independent
because the changed four-seed list was a basis.

After the same identifications $0=4$ and $6=7$, the eight remaining
crossing triples are identical, label for label, in the following pairs:
\[
\begin{array}{c|c|c|c}
\text{new knot}&\text{earlier reduced system}&W_0&
   \text{canonical third root}\\\hline
13a_{3264}&13a_{373}&u_{10}&u_8\\
13a_{3266}&13a_{375}&u_9&u_9\\
13a_{3275}&13a_{391}&u_9&u_9\\
13a_{3276}&13a_{392}&u_9&u_{12}.
\end{array}
\]
This is a finite incidence identity on the retained literal diagrams;
the checker reconstructs all four pairs directly, with no diagram
moves or assumed mutation invariance. By Lemma~\ref{lem:reduced},
all reduced normal directions span a space of dimension at most two.
But they include the independent three vectors $U,V,W_0$, a
contradiction. This proves the remaining cases.
\end{proof}

\section{Eight further alternating knots with six-sheeted actions}

Put
\[
\mathcal A_8=\{13a_{304},13a_{362},13a_{623},13a_{629},
13a_{3232},13a_{3259},13a_{3374},13a_{3375}\},\qquad
\mathcal K_{16}=\mathcal K_8\cup\mathcal A_8.
\]
The literal diagrams, canonical arc convention, original seed traces
and named-census transfers for $\mathcal A_8$ are retained in
\texttt{a6\_affine\_inputs.json}. The next proof uses their crossing
incidences, rather than their finite permutation or affine labels.

\begin{theorem}\label{thm:a6eight}
For every knot in $\mathcal A_8$, the four retained seed normal
directions are dependent in every common-form nonisotropic orthogonal
reflection representation over a field of characteristic different
from two. In particular, all meridian normal directions span a space
of dimension at most three. The ambient symmetric form may be singular.
\end{theorem}
\begin{proof}
Suppose the four original seed normals are independent. Write
$A=u_0$, $E=u_1$. The following table specifies a normal basis obtained
by changing only the last two original seeds. In each row the first
knot uses local type I below, and the second uses type II.
\[
\begin{array}{c|c|c|c|c}
 K&\text{original seeds}&\text{changed seeds}&C&\text{paired arc}\\\hline
13a_{304},13a_{3232}&(0,1,5,9)&(0,1,6,9)&u_6&7\\
13a_{362},13a_{3259}&(0,1,5,11)&(0,1,6,11)&u_6&7\\
13a_{623},13a_{3374}&(0,1,5,8)&(0,1,5,8)&u_8&9\\
13a_{629},13a_{3375}&(0,1,5,7)&(0,1,5,8)&u_8&9
\end{array}
\]
For the first two rows put $P=u_5$ and let $Q$ be the last
original seed. The literal crossing gives $u_6=R_QP=P-B(P,Q)Q$,
so replacing $P$ by $u_6$, while keeping $Q$, has determinant one.
In the last row $u_8=R_{u_5}u_7$ gives the other triangular change
of determinant one. The third row requires no change. These statements
hold before imposing any Gram restriction. The reduced directions
therefore include three independent normals: respectively
$u_0,u_6,u_9$; $u_0,u_6,u_{11}$; and $u_0,u_5,u_8$ in each
of the last two rows.

Let $j$ be the arc of $C$, and let $j'$ be the paired arc in the
last column. The five local crossing relations are
\[
\begin{array}{c|l}
\mathrm{I}&(0,1,2),(j,2,3),(3,0,1),(1,3,4),(2,j,j')\\
\mathrm{II}&(j,1,2),(0,2,3),(3,0,1),(1,3,4),(2,j,j').
\end{array}
\]
They are precisely the crossings involving one of arcs $1,2,3$.
Set $a=B(A,E)$, $b=B(A,C)$ and $c=B(E,C)$.
In type I choose
\[
 u_2=E-aA,\qquad u_3=E-aA-tC,\qquad t=c-ab,
 \qquad H=B(u_3,A)=-a-bt.
\]
The closing relation $R_{u_3}A=\varepsilon E$, $\varepsilon=\pm1$,
has coefficients
\[
                   1+aH=0,\qquad -H=\varepsilon,
                   \qquad tH=0.
\]
Thus $H\ne0$, $t=0$ and $a=\varepsilon$.
No scalar is cancelled before it is known nonzero.
Consequently $u_3=u_2$, $u_4=R_Eu_3=-aA$, and
$u_{j'}=R_{u_2}C=C$.

In type II choose instead
\[
 u_2=E-cC,\qquad u_3=E-cC-tA,\qquad t=a-bc,
 \qquad H=B(u_3,A)=-t.
\]
The same closing relation gives
\[
                   1+tH=0,\qquad -H=\varepsilon,
                   \qquad cH=0.
\]
It follows that $c=0$ and $t=a=\varepsilon$.
Again $u_4=-aA$ and $u_{j'}=C$.
Both types therefore identify reflection lines $0=4$ and $j=j'$.
The calculations do not require a nonzero Gram determinant.

Remove the five local crossings and fold these paired lines.
There remain eight crossing relations on eight normal directions.
Their exact relabelings to the already proved systems of
Lemma~\ref{lem:reduced} are listed below. A row gives the new arc
list in its displayed order, followed by the corresponding old labels;
the same row applies to both knots named in the first column.
\[
\begin{array}{c|c|c|c}
 K&\text{new labels}&\text{old labels}&\text{old system}\\\hline
304,3232&(0,5,6,8,9,10,11,12)&(9,10,11,12,0,5,6,8)&13a_{375}\\
362,3259&(0,5,6,8,9,10,11,12)&(10,11,12,0,5,6,8,9)&13a_{375}\\
623,3374&(0,5,6,7,8,10,11,12)&(8,6,5,0,12,11,10,9)&13a_{373}\\
629,3375&(0,5,6,7,8,10,11,12)&(5,0,12,11,10,9,8,6)&13a_{373}
\end{array}
\]
In the first column all indices denote $13a$ knots. The checker
reconstructs every crossing set from its literal PD and checks these
bijections, allowing only reversal of the two under-arcs.
By Lemma~\ref{lem:reduced}, all eight reduced normal directions span
a space of dimension at most two. They include the three independent
normals identified above, a contradiction.

The four retained seeds generate the whole knot group by their
complete propagation trace. Their normal span is invariant under the
seed reflections, and hence contains all conjugate meridian normals.
Its dimension at most three therefore gives the final assertion.
\end{proof}

\section{Coxeter consequences}

\begin{lemma}\label{lem:span}
Let $(W,S)$ be a Coxeter system. In its canonical real representation,
if $W$ is generated by reflections with normal vectors
$u_1,\ldots,u_d$, then these vectors span the full representation
space. In particular, its dimension, or Coxeter rank, is at most $d$.
If the rank is $d$, those normal vectors are a basis even when the
canonical form is singular.
\end{lemma}
\begin{proof}
For clarity, only existence of the canonical representation is needed
from the classical construction \cite[Section 6.12]{davis}.
Take a vector space with basis $e_s$, $s\in S$, and symmetric form
\[
 B(e_s,e_s)=2,\qquad
 B(e_s,e_t)=-2\cos(\pi/m_{st})
\]
for finite $m_{st}$, and take $B(e_s,e_t)=-2$ when $m_{st}=\infty$.
The maps $R_{e_s}$ satisfy the Coxeter relations. For a finite label,
the form on the two-normal plane is nondegenerate and positive
definite; splitting off this plane, the product is a rotation of
order $m_{st}$ and the identity on its orthogonal complement.
This argument remains valid if the full form has a radical.
There is no relation to check for an infinite label.
Thus the defining presentation gives a homomorphism from $W$ to
these linear maps, without using a faithfulness theorem.
Every Coxeter reflection maps to $R_{we_s}$ for a norm-two vector.

Put $E=\langle u_1,\ldots,u_d\rangle$. Each generating reflection
preserves $E$ and displaces every ambient vector by an element of
$E$. Induction on words gives the same displacement property for
the entire image of $W$. In particular,
$(R_{e_s}-I)e_s=-2e_s$ belongs to $E$ for every $s\in S$.
All standard normal basis vectors therefore lie in $E$.
This proves the claim even if $S$ was initially infinite.
\end{proof}

\begin{corollary}\label{cor:coxeter}
For each $K\in\mathcal K_{16}$, every
epimorphism from the group of $K$ to a Coxeter group
which takes meridians to Coxeter reflections has target Coxeter rank
at most three. In particular, no finite or infinite Coxeter reflection
quotient supplies the sharp bound $\mr(K)=4$.
\end{corollary}
\begin{proof}
For each knot the original four seeds generate its group by the
retained complete propagation trace, so their images generate any
epimorphic target. Lemma~\ref{lem:span} says
that their normal vectors span the canonical representation space.
Theorems~\ref{thm:alternating8} and~\ref{thm:a6eight} say that their
span has dimension at most three. Neither result requires an invertible Gram matrix or a
faithful representation.
\end{proof}

\begin{corollary}\label{cor:separation}
Each $K\in\mathcal K_{16}$ has $\mr(K)=b(K)=4$, while every Coxeter
reflection quotient has rank at
most three.
\end{corollary}
\begin{proof}
For $K\in\mathcal K_8$, the retained cover group is $C_3$, with
moving index two on seven sheets, so the arbitrary-cycle cover bound
is $\lceil(6+1)/2\rceil=4$. For $K\in\mathcal A_8$, the retained
ternary block matrices on eight dimensions have common fixed space
of dimension one and meridian displacement rank two; their fixed-space
bound is $\lceil7/2\rceil=4$, as established in the companion
\texttt{A6\_AFFINE\_SHARP\_FOUR.tex}, Theorem~2. Every literal
diagram also has a
four-seed bridge upper certificate. Corollary~\ref{cor:coxeter}
proves the asserted limitation of all Coxeter reflection quotients.
\end{proof}

\begin{corollary}\label{cor:bbk}
Each $K\in\mathcal K_{16}$ lies outside every
sharp-four construction in
the work of Baader, Blair and Kjuchukova \cite{bbk} whose proof supplies
a rank-four
Coxeter reflection quotient, including the rank-preserving rational
tangle replacements following the proof of their Theorem 1 in
Section 4.
\end{corollary}
\begin{proof}
Their Section 2, Proposition 1 and Section 4, Lemma 1 supply the
meridian-preserving Coxeter epimorphism furnishing the lower bound.
The stated rational replacements preserve that lower bound and the
matching upper bound, so a sharp-four application retains a rank-four
reflection quotient. Corollary~\ref{cor:coxeter} rules it out.
\end{proof}

This conclusion concerns the exact quotient-producing constructions.
It does not assert that arbitrary rational tangle substitutions
preserve rank or bridge number, nor classify all families called
generalized Montesinos.

\section{Exact algebra checks and assurance}

The eight six-sheet-action controls verify both raw local Gram
calculations, division-free reductions for both closing signs, all
triangular basis changes and the eight explicit reduced-system
relabelings. Their literal PDs and seed traces come from
\texttt{a6\_affine\_inputs.json}; the optional integer-15 certificates
are not needed for these transfers.

The final four alternating controls separately verify their raw
unit-determinant basis changes, alternate first-loop formulas,
normalized four-variable Gram form and the eight-crossing incidence
identity with the earlier cases. Their three-root propagation reaches
every reduced direction, including the independently known third
normal; no coefficient is assumed invertible to change that third root.

The sibling controls separately verify their original and changed seed
systems, all local crossings and reduced traces, and every coefficient
identity used in the three short contradictions. The $13a_{392}$
contradiction also has four generated integer identities with constant
one, avoiding division in its auxiliary symbols.

The standard-library checker first reconstructs the canonical arc
numbering, trace and unused crossings from the retained PD. It verifies
the raw six-variable first-loop formulas and all four coefficients
of $R_{u_0}(e_0)$ before imposing any reduction, and separately
checks both choices of the seed-root sign normalization.
It verifies the full normalized four-coordinate trace for
\[
 B=\begin{pmatrix}
 2&1&c&p\\1&2&0&p+b\\c&0&2&d\\p&p+b&d&2
 \end{pmatrix}
\]
and checks that every relevant direction equals its three-coordinate
formula in $U=e_0-e_1,V=e_2,W=e_3$. Thus the free parameter $p$
and singular forms have not been lost in the reduction.

The short contradiction also has a division-free integer check.
Treat $y,h,v,w,p,q$ as independent symbols and, for each pair
$\alpha,\beta\in\{1,-1\}$, form
\[
\begin{array}{lll}
 F_0=h-1+y^2,&F_1=p(y-w),&F_2=1+pv,\\
 F_3=-ph-\beta,&F_4=-1-q(1-wy),&F_5=-qwv-\alpha.
\end{array}
\]
Define successively
\[
 D=w-y,\quad E=v-\beta h,\quad A=-1-qh,\quad B_0=\beta y-\alpha.
\]
The following exact identities express each in the ideal $(F_0,\ldots,F_5)$:
\begin{align*}
 D&=vF_1+DF_2,\\
 E&=-\beta hF_2-\beta vF_3,\\
 A&=F_4-qF_0-qyD,\\
 B_0&=F_5-\beta yA+qvD+qyE,\\
 h&=F_0-(\beta y+\alpha)B_0,\\
 1&=-A-qh.
\end{align*}
Substituting these identities recursively gives explicit integer
multipliers with sum one. The checker verifies all four sign cases
by exact sparse polynomial multiplication. This part of the replay
uses no characteristic-zero division and supports the stated field
scope of the written proof.

As an optional independent discovery check, eight reduced sign-branch
certificates in three variables $c,b,d$, with Gram
$\left(\begin{smallmatrix}2&c&-b\\c&2&d\\-b&d&2\end{smallmatrix}\right)$,
satisfy integer identities $\sum_i H_i f_i=15$.
There are 285 multiplier monomials per branch, of degree at most nine.
The checker reconstructs all equations and multiplies the retained
integer coefficients. Those identities are not used at
characteristics three or five and are unnecessary for the short proof.
No Gr\"obner-basis implementation is trusted by the replay.
A corrupted multiplier is rejected. A singular rank-one Gram form on
four independent basis roots is also checked as a positive control:
its reflections still have independent displacement directions,
showing why Gram invertibility would be an unjustified hypothesis.

These checks support exact finite algebra and the retained diagram.
The uniform field and Coxeter statements require the written arguments;
no formal verification or external specialist review is asserted.
The separate priority audit must still compare other methods and prior
results before declaring the named equality globally new.

\begin{thebibliography}{9}
\bibitem{davis} M. W. Davis, \emph{The Geometry and Topology of
Coxeter Groups}, second edition, Springer Monographs in Mathematics
(2025), Section 6.12, printed pp.~108--111.
\url{https://people.math.osu.edu/davis.12/second_edition.pdf}.
\url{https://doi.org/10.1007/978-3-031-91303-7}.
\bibitem{bbk} S. Baader, R. Blair and A. Kjuchukova,
Coxeter groups and meridional rank of links,
Math. Ann. 379 (2021), 1533--1551, Sections 2 and 4.
\url{https://doi.org/10.1007/s00208-020-02124-z}.
\end{thebibliography}
\end{document}
